\subsection{Structural reductions for three objects and twelve arrows}

\begin{lemma}[A singleton cross row or column is a join]
\label{C:singleton-cross-join}
Let a category have objects $0<1<2$, with cross arrows only
in the increasing directions. If both $C(0,1)$ and $C(0,2)$
are singletons, then
\[
 C=\operatorname{End}(0)\star C|_{\{1,2\}},
\]
where $\star$ is the categorical join. Dually, singleton
$C(0,2)$ and $C(1,2)$ give
$C=C|_{\{0,1\}}\star\operatorname{End}(2)$.
In either case every positive-degree integral cohomology
product of $BC$ vanishes.
\end{lemma}
\begin{proof}
The categorical join retains the two full subcategories and
has exactly one arrow from every object of the first to every
object of the second. Its cross compositions are forced.
These are exactly the given hom-sets and compositions, including
the actions of an arbitrary extreme endomorphism monoid.
The nerve of a categorical join realizes to the topological
join of the two classifying spaces. Both spaces are nonempty
and connected here, so their join is homotopy equivalent to
the suspension of their smash product. Positive-degree products
on a suspension vanish. The opposite case is identical.
\end{proof}

\begin{lemma}[A rational consequence of the all-act obstruction]
\label{C:eight-bounded-cups}
Consider a triangular two-object category with one endomorphism
monoid $M$ of order at most eight and another monoid $N$ whose
trivial module is projective over the rational algebra on the
side needed for its coinvariants. If its rational homology
vanishes in all sufficiently high degrees, every product of
two degree-two rational cohomology classes is zero.
\end{lemma}
\begin{proof}
Projectivity of the trivial module makes $BN$ rationally
acyclic and its coinvariant functor exact. If
$H^2(BM;\mathbb Q)=0$, the relative-suspension argument of
Lemma~\ref{C:small-triangular-vertices} proves the assertion.
Otherwise the complete proof of
Proposition~\ref{C:eight-act-cones}, applied to the exact
coinvariant augmentation kernel in
Lemma~\ref{C:act-cone-Tor}, says that nonzero fourth rational
homology would persist in arbitrarily high degrees. Boundedness
therefore forces $H_4=0$, which also forces every displayed
degree-two product to be zero. This argument uses the proved
Tor criterion, rather than an inference from the homotopy-type
statement alone.
\end{proof}

\begin{lemma}[A singleton leg in a star]
\label{C:star-singleton-join}
Let a three-object category have cross arrows only on two
legs of a star, both directed toward the center or both away
from it. If one cross-hom set is a singleton and the opposite
leaf has rationally acyclic classifying space, every positive
degree rational cohomology product of the category vanishes.
The endomorphism monoid at the singleton leaf may be arbitrary.
\end{lemma}
\begin{proof}
Write $X$ for the center, $Y$ for the singleton leaf, and $Z$
for the other leaf, using the same letters for their
classifying spaces. The singleton collage is $J=X*Y$,
and the inclusion $X\to J$ is null-homotopic. Write $D$ for
the other collage and $E$ for its two-sided edge bar. Since
there are no composable pairs of cross arrows, the classifying
space of the star is exactly the homotopy pushout of the two
collages along $X$. Nullity of the first inclusion gives
\[
 BC\simeq J\vee\operatorname{hocofib}(X\to D)
       \simeq J\vee\operatorname{hocofib}(E\to Z).
\]
The join of the connected based spaces $X,Y$ is homotopy
equivalent to a suspension. Since $Z$ is rationally acyclic,
its map to a point induces a rational homology equivalence
from the second summand to the unreduced suspension of $E$.
These maps induce cohomology ring isomorphisms over $\mathbb Q$.
Both summands consequently have zero positive-degree products,
as does their wedge.
\end{proof}

\begin{lemma}[All three-object stars through twelve arrows]
\label{C:twelve-all-stars}
No three-object star category with at most twelve arrows has
classifying space homotopy equivalent to $\mathbb CP^2$.
\end{lemma}
\begin{proof}
Let $m,n,p$ be the three endomorphism orders and let $s,t\geq1$
be the leg cardinalities. Thus $m+n+p+s+t\leq12$.
Every monoid of order at most four has rationally acyclic
classifying space, by Theorem~\ref{E:all-six-types}.
In particular, a vertex with nonzero rational $H^2$ has
endomorphism order at least five. There can be at most one
such vertex, since two orders at least five, a third identity,
and two legs would already require thirteen arrows.

If all three vertex spaces have zero rational $H^2$, every
degree-two class lifts to cohomology relative to their union.
After collapsing that union, the graph-of-spaces description
is a wedge of suspensions of the two edge bars with disjoint
basepoints. Relative degree-two products are therefore zero,
and so are the absolute products. This excludes $\mathbb CP^2$.

Suppose first that the unique vertex with nonzero rational
$H^2$ is the center, of order $m\geq5$.
The leaves have order at most four and are rationally acyclic.
A singleton leg is handled by
Lemma~\ref{C:star-singleton-join}. Otherwise $s,t\geq2$ and
the budget gives $n+p\leq3$. Both leaf algebras are therefore
semisimple over $\mathbb Q$, since their monoids have order
one or two. Take their exact coinvariants in the two edge
bars. If the resulting right $\mathbb QM$-modules are $T_1,T_2$,
the high-degree part of the total star chain complex is the
mapping cone of
\[
 (T_1\oplus T_2)\otimes_{\mathbb QM}P_\bullet
 \longrightarrow\mathbb Q\otimes_{\mathbb QM}P_\bullet,
\]
where $P_\bullet$ is a projective resolution of the trivial
module and the map is the sum of the two augmentations.
The two leaf copies of $\mathbb Q$ affect only degrees zero
and one. Thus, for $q\geq2$,
\[
 H_q(BC;\mathbb Q)
 \cong\operatorname{Tor}_{q-1}^{\mathbb QM}(L,\mathbb Q),
 \qquad L=\ker(T_1\oplus T_2\to\mathbb Q).
\]
Here $m\leq8$ and $L$ is finite-dimensional. The same
all-module Tor criterion used in
Proposition~\ref{C:eight-act-cones} says that nonzero fourth
homology persists arbitrarily high. This contradicts the
homology of $\mathbb CP^2$.

Finally suppose the distinguished vertex is a leaf, of order
$m\geq5$. The center and the other leaf are rationally acyclic.
Collapsing the center therefore induces a rational cohomology
ring isomorphism and splits the star into a wedge of its two
leg cofibers. Each leg cofiber is also rationally equivalent,
by a cohomology ring isomorphism, to its original collage.
The collage between the two acyclic vertices has zero
positive-degree products. A nonzero degree-two square must
therefore occur in the collage containing the distinguished
leaf, and that collage has bounded rational homology if the
whole star does.

Let $s$ be this distinguished leg's cardinality and $n$ the
center's endomorphism order. If $s\leq2$,
Lemma~\ref{C:two-object-cup-pruning} gives zero products.
If $s\geq3$, the other leaf and other leg contribute at least
two arrows, so
\[
 m+n\leq12-s-2\leq7.
\]
Hence $n\leq2$, its rational algebra is semisimple, and the
distinguished leaf has order at most eight.
Lemma~\ref{C:eight-bounded-cups} again gives zero degree-two
products. This is the final contradiction.
\end{proof}

\begin{proposition}[Only the linear three-object condensation remains]
\label{C:twelve-three-condensation}
Let $C$ be a three-object category with at most twelve arrows
whose classifying space is homotopy equivalent to $\mathbb CP^2$.
After labelling its objects, it must have cross arrows exactly
in the three directions
\[
 0\longrightarrow1,\qquad1\longrightarrow2,\qquad0\longrightarrow2,
\]
with all three displayed hom-sets nonempty and no reverse arrows.
\end{proposition}
\begin{proof}
Any object that is a retract of another may be deleted without
changing the classifying space. That would leave at most eleven
arrows, contradicting Proposition~\ref{C:CP2-lower12}.
We may therefore assume there are no such retracts.

If all three objects form one strongly connected component,
apply Lemma~\ref{C:finite-category-Rees}. None of their
identities belongs to its minimal ideal, by the retract
assertion. If neither all row numbers nor all column numbers
are one, both total numbers are at least four, so the ideal
alone contains at least sixteen arrows. Otherwise the
classifying space has finite-group type and is rationally
acyclic. Both alternatives exclude $\mathbb CP^2$.

Suppose the strongly connected components have sizes two
and one. The two-object block has at least six arrows:
both endomorphism monoids have order at least two, since a
trivial one would make that object a retract, and there are
arrows in both directions. Connectedness supplies a cross
arrow between this block and the singleton. Composing with
arrows inside the block supplies cross arrows for both block
objects, so there are at least two such arrows.
The singleton's endomorphism order is therefore at most four,
and its classifying space is rationally acyclic.

Conversely, the two-object block has at most nine arrows.
By Lemma~\ref{C:finite-category-Rees}, a block without a
retract that is not of finite-group type needs total row and
column numbers at least three, nine minimal-ideal arrows,
and its two identities outside that ideal. It would therefore
need at least eleven arrows. Our block is consequently of
finite-group type and is rationally acyclic as well.

The mixed-nerve bar decomposition now writes $BC$ as a
homotopy pushout of the two diagonal block spaces over their
cross-bimodule bar. Replacing the two rationally acyclic
vertex spaces by points gives a rational homology equivalence
to a suspension, by the homotopy-pushout exact sequence.
It induces a rational cohomology ring isomorphism, so all
positive products vanish, again a contradiction.

The three strongly connected components are therefore
singletons. Their connected condensation poset is either a
two-leg star or a linear order. The star is excluded by
Lemma~\ref{C:twelve-all-stars}. In a linear order, composition
of the two adjacent nonempty hom-sets makes the third
displayed hom-set nonempty. This is exactly the asserted
remaining configuration.
\end{proof}
